\documentclass[../../../include/open-logic-section]{subfiles}

\begin{document}

\olfileid[fr]{sfr}{infinite}{card-sb}

\olsection{Annexe : une démonstration de Schröder--Bernstein}

Avant de quitter la théorie naïve des ensembles, il nous reste une
dernière preuve naïve, mais élaborée, à examiner. Il s'agit d'une preuve
du théorème de Schröder--Bernstein
(\olref[sfr][siz][sb]{thm:schroder-bernstein}) : si $\cardle{A}{B}$ et
$\cardle{B}{A}$, alors $\cardeq{A}{B}$. Autrement dit, étant données
des !!{injection}s $f\colon A\to B$ et $g\colon B\to A$, il existe
!!a{bijection} $h\colon A\to B$.

Dans ce chapitre, nous avons suivi l'idée de \emph{clôture} de Dedekind.
Dedekind a justement donné une belle preuve de Schröder--Bernstein
fondée sur cette notion ; nous allons la présenter ici. Notre preuve
suit de près celle de \citet[pp.~157--8]{Potter2004}, auquel on pourra
se reporter pour une présentation un peu différente, mais semblable
pour l'essentiel. Une brève recherche sur le Web permet aussi de constater
que ce théorème, tout comme l'irrationalité de $\sqrt{2}$, admet de
\emph{nombreuses} preuves différentes et intéressantes.

Avec une notation analogue à celle de la
\olref[sfr][infinite][dedekind]{Closure}, posons, pour une application
$f\colon U\to U$ et une partie $B\subseteq U$,
\[
\Closureofunder{f}{B} = \bigcap \Setabs{X\subseteq U}{B \subseteq X
\text{ et }X\text{ est stable par }f}.
\]
Ainsi définie, $\Closureofunder{f}{B}$ est la plus petite partie de $U$
stable par $f$ qui contient $B$, au sens suivant.\footnote{Note de
l'édition française : comme pour la clôture d'un élément, on fixe
l'ensemble ambiant et on suppose que l'application le préserve.
La famille de parties à intersecter contient $U$ ; elle est donc
non vide, y compris lorsque $B$ est vide.}

\begin{lem}\ollabel{Closureprops}
	Pour toute application $f\colon U\to U$ et toute partie $B\subseteq U$ :
	\begin{enumerate}
		\item\ollabel{Closurehaselem} $B \subseteq \Closureofunder{f}{B}$ ;
		\item\ollabel{Closureclosed} $\Closureofunder{f}{B}$ est stable par $f$ ;
		\item\ollabel{Closuresmallest} si $X\subseteq U$ est stable par $f$
		et si $B\subseteq X$, alors $\Closureofunder{f}{B}\subseteq X$.
	\end{enumerate}
\end{lem}

\begin{proof}
La preuve est exactement la même que celle du
\olref[sfr][infinite][dedekind]{closureproperties}.
\end{proof}

Un dernier résultat nous permettra d'obtenir le théorème de
Schröder--Bernstein.

\begin{prop}\ollabel{sbhelper}
Si $A \subseteq B \subseteq C$ et $A \approx C$, alors
$\cardeq{\cardeq{A}{B}}{C}$.
\end{prop}

\begin{proof}
Étant donnée !!a{bijection} $f\colon C\to A$, considérons-la aussi
comme une application de $C$ dans $C$, puisque $A\subseteq C$.
Posons $F=\Closureofunder{f}{C\setminus B}$ et définissons une
fonction $g$ de domaine $C$ par
\[
	g(x) =
	\begin{cases}
			f(x) &\text{si $x \in F$},\\
			x &\text{sinon}.
		\end{cases}
\]
Nous allons montrer que $g$ est !!a{bijection} de $C$ sur $B$.
Il s'ensuivra que $\comp{f^{-1}}{g}\colon A\to B$ est
!!a{bijection}, ce qui achèvera la preuve.

Montrons d'abord que, si $x\in F$ et $y\notin F$, alors
$g(x)\neq g(y)$. Supposons le contraire pour obtenir une contradiction :
on aurait $y=g(y)=g(x)=f(x)$. Or $x\in F$ et $F$ est stable par $f$,
d'après le \olref{Closureprops}. On aurait donc $y=f(x)\in F$,
ce qui est contradictoire.

Supposons maintenant que $g(x)=g(y)$. D'après ce qui précède,
$x\in F$ si et seulement si $y\in F$. Si $x,y\in F$, alors
$f(x)=g(x)=g(y)=f(y)$ ; ainsi, $x=y$, puisque $f$ est !!a{bijection}.
Si $x,y\notin F$, alors $x=g(x)=g(y)=y$. Dans les deux cas, $x=y$ :
$g$ est donc !!a{injection}.

Il reste à montrer que $\ran{g}=B$. Vérifions d'abord que
$\ran{g}\subseteq B$ : si $z\in F$, alors $g(z)=f(z)\in A\subseteq B$ ;
si $z\in C\setminus F$, alors $z\notin C\setminus B$, puisque
$C\setminus B\subseteq F$, donc $g(z)=z\in B$.\footnote{Note de
l'édition française : cette inclusion, implicite dans la source, est
nécessaire pour établir que $g$ a bien ses valeurs dans $B$.}
Fixons maintenant $x\in B\subseteq C$. Si $x\notin F$, alors $g(x)=x$.
Si $x\in F$, il existe $y\in F$ tel que $x=f(y)$ : sinon
$F\setminus\{x\}$ serait stable par $f$ et contiendrait $C\setminus B$.
En effet, l'image de chacun de ses éléments serait dans $F$ par stabilité,
et serait différente de $x$ par l'absence supposée d'antécédent dans $F$ ;
de plus, retirer $x\in B$ ne retire aucun élément de $C\setminus B$.
Cela contredirait la minimalité donnée par le \olref{Closureprops}.
Nous avons donc $g(y)=f(y)=x$, ce qui prouve l'inclusion réciproque.
\end{proof}

Voici enfin la preuve du résultat principal. Rappelons que, pour une
fonction $h$ et un ensemble $D$ inclus dans son domaine, on définit
$\funimage{h}{D}=\Setabs{h(x)}{x\in D}$.

\begin{proof}[Preuve du théorème de Schröder--Bernstein]
Soient $f\colon A\to B$ et $g\colon B\to A$ des !!{injection}s.
Puisque $\funimage{f}{A}\subseteq B$, nous avons
$\funimage{g}{\funimage{f}{A}}\subseteq g[B]\subseteq A$.
De plus, $\comp{f}{g}\colon A\to\funimage{g}{\funimage{f}{A}}$
est une !!{injection}, puisque $g$ et $f$ sont toutes deux injectives ;
c'est même !!a{bijection}, par la définition de son codomaine.
Ainsi, $\cardeq{\funimage{g}{\funimage{f}{A}}}{A}$, et la
\olref{sbhelper} fournit donc !!a{bijection}
$h\colon A\to\funimage{g}{B}$. En outre, $g^{-1}$ est !!a{bijection}
de $\funimage{g}{B}$ sur $B$. Par conséquent,
$\comp{h}{g^{-1}}\colon A\to B$ est !!a{bijection}.
\end{proof}

\end{document}
